%%%%%%%%%%%%%%%%%%%%%%%%%%%%%%%%%%%%%%%%%%%%%%%%%%
%%                                              %%
%%         Capítulo 1 - Introdução              %%
%%                                              %%
%%%%%%%%%%%%%%%%%%%%%%%%%%%%%%%%%%%%%%%%%%%%%%%%%%

\chapter{Introdução}
\label{cap:introducao}

\section{Contextualização}

A transição de sistemas monolíticos para arquiteturas de microsserviços --- processo denominado \textit{microservitização} --- tem sido amplamente impulsionada pela necessidade de alinhar decisões técnicas de design com a criação de valor de negócio \cite{hassan2020microservice}. Ao isolar capacidades funcionais em serviços autônomos e desacoplados, busca-se aumentar a escalabilidade, a facilidade de manutenção e a manutenibilidade dos sistemas distribuídos em nuvem \cite{verarivera2021granularity}.

No entanto, definir o tamanho e a dimensão ideal de cada serviço (o problema de granularidade) permanece como um dos maiores desafios em engenharia de software \cite{hassan2020microservice, verarivera2021granularity}. A decomposição excessiva (serviços muito pequenos ou \textit{nanoserviços}) pode introduzir complexidade operacional desnecessária e sobrecarga (\textit{overhead}) de rede, enquanto serviços demasiadamente agregados anulam os benefícios de isolamento e escalabilidade independente \cite{mendonca2024serviceweaver, rego2025granularity}.

\section{Motivação}
\label{sec:motivacao}

Decisões de granularidade tomadas no momento do projeto tendem a se perpetuar ao longo da vida do sistema: reorganizar serviços após a implantação --- seja dividindo um serviço sobrecarregado, seja agregando serviços excessivamente fragmentados --- envolve retrabalho significativo, coordenação entre equipes e riscos de migração em produção \cite{hassan2020microservice}. Ainda assim, a escolha é frequentemente realizada sem evidências objetivas, por heurísticas de moda ou intuição, e seus erros se manifestam de forma custosa: sobrecarga de comunicação que degrada o desempenho, ou perda de escalabilidade que exige provisionamento desproporcional de infraestrutura \cite{gan2019deathstarbench, mendonca2024serviceweaver}.

Os estudos empíricos recentes deram passos importantes ao quantificar os efeitos da granularidade sobre o desempenho \cite{mendonca2024serviceweaver} e sobre o desempenho e a resiliência \cite{rego2025granularity}. Contudo, esses efeitos são reportados como atributos de qualidade isolados, sem uma consolidação econômica que indique, para um perfil de carga concreto, qual configuração oferece o melhor equilíbrio entre o que se ganha e o que se paga --- lacuna que motiva diretamente a presente proposta.

Completa essa motivação a viabilidade prática do estudo: o arcabouço \textit{Service Weaver} \cite{ghemawat2023serviceweaver} torna a granularidade uma configuração de implantação, permitindo conduzir experimentos controlados com uma aplicação de referência realista --- a \textit{Online Boutique} --- a um custo computacional e de desenvolvimento acessível a um trabalho de conclusão de curso. O tema interessa, ainda, à formação em Sistemas de Informação, por situar a análise econômica como critério explícito de decisão arquitetural.

\section{Formulação do Problema}

Os métodos tradicionais de avaliação econômica e de atributos de qualidade em arquitetura de software, como o \textit{Cost Benefit Analysis Method} (CBAM) \cite{asundi2001cbam}, foram projetados para serem aplicados de forma estática, predominantemente na fase de design. Contudo, a granularidade ideal de um microsserviço não é uma propriedade fixa, pois depende diretamente das condições operacionais do sistema (como perfis de tráfego, variação de latência e consumo de recursos) \cite{hassan2020microservice, hassan2017ambients}.

Nesse contexto, há carência de evidências empíricas que quantifiquem, para cenários de carga específicos, o equilíbrio entre duas dimensões cruciais na escolha da granularidade \cite{hassan2020microservice}:
\begin{enumerate}
    \item \textbf{O valor agregado (benefício):} ganhos de escalabilidade e atendimento aos acordos de nível de serviço (SLAs) obtidos ao desacoplar componentes.
    \item \textbf{O custo incorrido:} impactos técnicos e financeiros decorrentes da sobrecarga de comunicação entre processos e máquinas virtuais.
\end{enumerate}

O surgimento recente de arcabouços como o \textit{Service Weaver} \cite{ghemawat2023serviceweaver} oferece uma oportunidade única para tratar esse problema. Ao permitir que a aplicação seja desenvolvida como um monólito modular e compilada em um único binário, o \textit{Service Weaver} possibilita alterar a granularidade física de implantação (co-localizando ou desacoplando componentes via arquivo de configuração TOML) sem alterações no código-fonte de negócio \cite{mendonca2024serviceweaver}. Isso viabiliza um estudo experimental comparativo: implantar uma mesma aplicação em diferentes granularidades e medir, sob carga controlada, os custos e benefícios de cada configuração.

\section{Objetivos}

\subsection{Objetivo Geral}

Avaliar empiricamente a relação custo-benefício de diferentes granularidades de implantação de microsserviços, utilizando o arcabouço \textit{Service Weaver} e a aplicação \textit{Online Boutique} sob cenários específicos de carga, de modo a identificar a configuração que oferece o melhor equilíbrio entre escalabilidade e custo.

\subsection{Objetivos Específicos}

\begin{itemize}
    \item Definir um conjunto de cenários de granularidade (topologias de implantação) para a aplicação \textit{Online Boutique}, variando a co-localização e o desacoplamento de componentes por meio da configuração TOML do \textit{Service Weaver}, tomando o CBAM como referencial de análise custo-benefício.
    \item Executar experimentos controlados com o gerador de carga \textit{Locust}, simulando perfis de carga específicos (usuários concorrentes e picos súbitos de acesso) sobre cada cenário definido.
    \item Analisar comparativamente as métricas coletadas (latência média e p95, vazão e número de chamadas remotas) para identificar a granularidade com melhor equilíbrio custo-benefício em cada perfil de carga.
\end{itemize}
